\section{Memory hierarchy}
\label{sec:memory}

\subsection{Levels}

Dependent chases place the levels at: L1 data 23.4~ns (34 cycles; the
read-only \texttt{\_\_ldg} path measures the same 34, one physical
array); L2 161.5~ns (235 cycles), flat from 128~KiB to 5~MiB with the
cliff binding between 5 and 8~MiB against the 6~MB L2; DRAM 299.9~ns.
Bandwidths: DRAM streams at 608~GB/s read (97.5\% of the P2 ceiling, and
within 0.3\% of an independent estimate from the sector experiment
below), 474 write, 507 copy. The L2's own service rate is lower than the
DRAM streaming rate: a 4~MiB footprint read entirely from L2 via
\texttt{.cg} sustains 382~GB/s (the request-serving path saturates
before the line-fill path), while the same footprint under the default
policy reads at 1.11~TB/s because it nearly fits the \emph{aggregate} L1
capacity (72 SMs $\times$ 64~KiB). Shared memory and L1 share one 64-byte
per-cycle-per-SM datapath (both measure the same ceiling by two load
widths), and the 96~KiB carveout moves the L1 chase cliff exactly as the
split predicts. One hazard row: \texttt{float4} shared-memory access
compiles to paired \texttt{LDS.64} whose 16-byte stride self-conflicts
two-way, halving bandwidth --- tiles built on \texttt{float4} pay it.

\subsection{TLBs, constant path, instruction caches}

With \texttt{.cg} loads (the virtually indexed L1 otherwise hides TLB
behaviour entirely), both TLB cliffs land on the T4 priors: $+7.0$~ns
past 32~MiB of reach and $+5.9$~ns past 8~GiB \citep[p.~38]{jia2019turing},
small penalties on a 160~ns L2 baseline. The constant path resolves into
three measured levels: a warp-uniform chain lowers to the uniform
datapath at 5.1 cycles, divergent-register chains hit the constant cache
at 26.1, and a 32-address indexed read shows the same per-step
\emph{latency} (serialisation costs throughput, not dependent-chain
time). Instruction-fetch cost stays at the single-warp issue floor
through 16~KiB loop bodies, softens at 32~KiB, and steps at 64 and
128~KiB --- bracketing the L0 and L1I capacities.

\subsection{The sector model, load policies, atomics}

The quantised-block stride experiment binds the row that motivated it: a
warp of byte loads at stride 18 (the q4\_0 block layout) requests 18
32-byte sectors, and measured useful bandwidth times sectors reproduces
the DRAM peak --- 610~GB/s of fetched sectors at stride 18 and 32, the
fetch-bound regime exactly where the model predicts it. Small strides
are request-rate-bound instead, and stride 128 additionally loses DRAM
row locality. Load-policy probes refute a common assumption:
\texttt{.cs} (evict-first) lines still enter the L1 and re-read at
23.4~ns --- the hint sets priority, not eviction --- while \texttt{.lu}
genuinely evicts and \texttt{.cg} never allocates (both 161~ns re-read).
Keeping weights resident against streamed data therefore means marking
the \emph{streamed} side \texttt{.cg}, not the weights \texttt{.cs}.
Shared-memory atomics reproduce the T4 contention deltas exactly
($+2/+4/+8/+16/+32$ cycles for 2- to 32-way) atop a constant
17-cycle offset attributable to the return-value chain this method times
and the prior's stall-counting method does not; global atomics sit at
${\sim}310$ cycles through the L2 with a real 2\% spread between the two
physical GPUs.
